\section{A square-lattice comparison for the harmonic frame}
\label{sec:harmonic-beta-square-lattice}

Put \(N=\ker(\lambda,\rho_2)\) in full profinite GT.
We use the full mixed genus-one action, its forgetting compatibility,
and its two low-genus descriptions from
\cite[Lemmas~16--17]{H17}. We also use the normalized harmonic identity
on all of \(N\), Proposition~\ref{prop:normalized-GF1-from-dyadic}, and
the signed congruence identity established above.
No compatibility with an arbitrary finite cover is assumed.
Here \(B_g^{\mathrm{ad}}\) denotes the normalized full adelic
Alexander unit in the completed group algebra
\(\widehat{\mathbf Z}[[\widehat{\mathbf Z}^{\,2}]]\):
\[
 B_g^{\mathrm{ad}}(X,Y)
       =1+(X-1)(Y-1)R_g^{\mathrm{ad}}(X,Y).
\]
Its value at \((i,-1)\) is defined by the continuous finite-character
map to \(\mathbf Z_\ell[i]\) at each coefficient prime \(\ell\).
At two, its augmentation branch also admits the usual dyadic formal
power-series description; the finite-character definition is used
at every odd prime. Throughout this section only, abbreviate
\(B_g^{\mathrm{ad}}\) by \(B_g\). Write
\[
 O_4=\mathbf Z[i],\qquad
 \nu_{4,g}=B_g(i,-1).
\]

\begin{theorem}\label{thm:harmonic-beta-square-lattice}
There is one integral identification of the harmonic quadratic order
with \(\widehat O_4\), valid for every \(g\in N\) and every coefficient
prime, under which the corrected harmonic frame \(W_g\) is
\(\nu_{4,g}\). Reversing the choice of primitive fourth root conjugates
all these values simultaneously.
Consequently, on \(D=N_{\rm rad}\), the tame harmonic square
test at an odd prime \(\ell\) is equivalent to
\[
 B_{g,\ell}(i,-1)^{(\ell-\chi_{-4}(\ell))/2}=1
\]
in the residue norm-one group. Either choice of the primitive
fourth root gives the same test.
\end{theorem}

\subsection{The finite translation group and the quotient action}

Let \(T=\mathbf R^2/\mathbf Z^2\) and mark all sixteen points
of \(T[4]\), with their individual labels. In the full mixed
mapping class group \(\Gamma=\Gamma_{1,[16]}\), the translations
form \(H=T[4]\simeq(\mathbf Z/4)^2\).
Let \(\alpha=\alpha_P\) be the mixed action of \(g\).
The finite-\(2\)-subgroup conjugacy theorem
\cite[Theorem~A(i); Theorem~6.2(i)]{CubicBZ}
puts \(\alpha(H)\) into a discrete finite subgroup by a pure
profinite conjugation. To obtain purity, apply the theorem
in the inverse image of the translation permutation group
\(H\); its finite quotient is abelian, so the discrete
translation part of an initial conjugator can be absorbed
without changing the labelled generators. The discrete group obtained in this
way has the same permutation on the sixteen labels as \(H\).
The homology covariance of \(\alpha\) makes its action on
closed-torus homology trivial. Nielsen realization therefore
realizes it by translations on a complex torus \(T'\)
\cite[Section~4]{CubicKerckhoff}, simply
transitive on the marked points. After choosing an origin,
the labelled translation group is an isomorphism
\(H\to T'[4]\).

We check the orientation of the marking that identifies these
two translation actions. After an oriented linear identification
of the tori, the labelled isomorphism is a matrix in
\(\mathrm{GL}_2(\mathbf Z/4)\). It lifts to
\(\mathrm{GL}_2(\mathbf Z)\): matrices of determinant one lift
from \(\mathrm{SL}_2(\mathbf Z/4)\), and the other determinant
class is represented by a reflection. Let its orientation
be \(\varepsilon\in\{1,-1\}\).
We spell out the point-pushing input that determines this sign.
Put \(\Pi=\pi_1(T\setminus T[4],*)\). Add a seventeenth marked
point and use the forgetting-compatible action on its stabilizer.
Its kernel is \(\widehat\Pi\), so this action supplies
\(\eta\in\operatorname{Aut}(\widehat\Pi)\) such that the
outer action of \(\alpha(h)\) is that of
\(\eta h\eta^{-1}\) for every \(h\in H\).
This is precisely the point-pushing construction and positive
puncture-loop assertion of \cite[Lemma~16(b),(c)]{H17}.
Let \(e\) be the initial pure profinite correction and let
\(r\) be the discrete labelled marking of orientation
\(\varepsilon\) above, so that
\(\alpha(h)=erhr^{-1}e^{-1}\).
The outer automorphism induced by \(\eta^{-1}er\) centralizes
the outer deck action of \(H\) on \(\widehat\Pi\), and it
carries each puncture loop to an \(\varepsilon\)-th power
up to \(\widehat\Pi\)-conjugacy. The signs here follow from
the positive loop assertion for \(\eta\), purity of \(e\),
and the orientation of \(r\).

Let
\[
 F=\pi_1\bigl((T\setminus T[4])/H\bigr),\qquad
 \widehat\Pi=\ker(\widehat F\longrightarrow H),
\]
where \(F\) is free of rank two and the quotient is
abelianization modulo four. Conjugation identifies
\(\widehat F\) with the full preimage of the outer deck
subgroup \(H\) under
\[
 \operatorname{Aut}(\widehat\Pi)
        \longrightarrow\operatorname{Out}(\widehat\Pi).
\]
Indeed its kernel is \(\operatorname{Inn}(\widehat\Pi)\),
its image on the outer quotient is exactly \(H\), and
its conjugation map is injective since the centralizer of
an open nonabelian free subgroup is trivial.
Choose an automorphism representing the outer class of
\(\eta^{-1}er\). Conjugation by it preserves this full
preimage and therefore defines
\(\Theta\in\operatorname{Aut}(\widehat F)\) acting as
that representative on \(\widehat\Pi\) and as the identity
on \(H\).
A positive boundary loop of the once-punctured quotient
torus is a puncture loop in \(\Pi\). Hence \(\Theta\)
sends it to a conjugate of its \(\varepsilon\)-th power.
In the class-two quotient it is the commutator of a free
basis, so \(\det\Theta^{\mathrm{ab}}=\varepsilon\).
Since \(\Theta\) acts identically modulo four, this
number is one modulo four. Thus \(\varepsilon=1\).
The resulting correction is pure and orientation preserving,
and we obtain
\begin{equation}
                  \alpha'|_H=\mathrm{id}_H.                 \label{eq:harmonic-beta-q1}
\end{equation}
This is the point-pushing and determinant argument of
\cite[Proposition~18, Step~1]{H17}, with the characteristic
four-division cover specified explicitly. It uses no
lifting-surjectivity assertion for a cyclic isogeny.

The discrete normalizer of \(H\) is
\(H\rtimes\mathrm{SL}_2(\mathbf Z)\). Its subgroup fixing
the origin is \(N_0=\mathrm{SL}_2(\mathbf Z)\), and forgetting
the other fifteen points identifies its completion with
\(\widehat\Gamma_{1,1}\). The finite-subgroup normalizer
theorem \cite[Theorem~A(iii); Theorem~6.2(iii)]{CubicBZ}
and \eqref{eq:harmonic-beta-q1} show that \(\alpha'\) preserves this completed
subgroup. Forgetting compatibility gives
\begin{equation}
     \alpha'|_{\widehat N_0}
        =\operatorname{Ad}(c_2)\Phi_g
              \quad(c_2\in\widehat\Gamma_{1,1}).             \label{eq:harmonic-beta-q2}
\end{equation}
Its central involution \(\iota=-1\) is fixed.
Moreover \(c_2\) is pure modulo two. Indeed \eqref{eq:harmonic-beta-q1} makes the
action on \(N_0\) trivial modulo four, and normalized
congruence triviality does the same for \(\Phi_g\).
Thus \(c_2\bmod2\) centralizes
\(\mathrm{SL}_2(\mathbf F_2)\), whose centre is trivial.

Take the centralizer of \(\iota\) in the subgroup fixing the
labelled origin. Its quotient by \(\langle\iota\rangle\) is
the mixed mapping class group \(\mathcal D\) of a sphere
with ten marked points: four branch points \(b_0,b_1,b_2,b_3\),
the images of \(T[2]\), and six other points, the images of
the pairs \(\{t,-t\}\subset T[4]\setminus T[2]\).
The group fixes \(b_0\) and permits permutations within the
other three branch labels and within the six pair labels.
The centralizer has commuting representatives: take the
unique extremal representative in a fixed punctured conformal
structure for which \(\iota\) is conformal. Composing on
either side with \(\iota\) gives extremal representatives
of the same class, so they agree. Existence and uniqueness
in this finite-type setting are the inputs
\cite[Sections~5 and~8, Theorems~1 and~6]{CubicGH}.
Every such mapping class lifts: after deleting all ten
points, the monodromy of the double cover is one at each
branch point and zero at the other six. The lifts differ
by \(\iota\). The unbranched fibre-isotopy theorem
\cite[Theorem~1]{CubicBH} makes descent independent of the
chosen commuting representative. This proves the discrete
quotient sequence using the unchanged part of that theorem.

The sequence has the required full profinite topology.
The upstairs and downstairs punctured-surface groups are
an open pair of nonabelian free groups. The completed
downstairs mixed mapping class group acts faithfully on the
completed downstairs surface group: Asada's theorem
\cite[introduction Theorem, (i)--(ii)]{CubicAsada} gives
the pure assertion, and a nontrivial permutation is visible
on peripheral conjugacy classes. Restriction to the open
upstairs group is faithful after quotienting by the deck
group, because an automorphism fixing an open free subgroup
pointwise is the identity. The deck involution itself is
detected by its permutation of the twelve nonfixed labels.
It follows that the completed centralizer embeds in the
completed ambient group. The finite-\(2\)-subgroup
centralizer theorem then identifies it with the full
profinite centralizer of \(\iota\).
Consequently \(\alpha'\) induces an automorphism
\(\delta_{10}\) of \(\widehat{\mathcal D}\).

On the pure ten-point group, genus-zero rigidity
\cite[Definition~2.7 and Corollary~2.8(i)]{CubicHMM} gives
\[
 [\delta_{10}]=(g',\sigma)
 \quad\text{in }\operatorname{Out}(\widehat\Gamma_{0,10})
       =\widehat{\mathrm{GT}}\times\mathfrak S_{10}.
\]
All sixteen labels were marked, and the correction in \eqref{eq:harmonic-beta-q1}
was pure. Thus the induced action on
\(\mathfrak S_3\times\mathfrak S_6\) is the identity.
The centralizer of this permutation subgroup in
\(\mathfrak S_{10}\) is trivial, so \(\sigma=1\).
To identify \(g'\), retain the four branch points and one
pair label. Upstairs this is forgetting all points except
the corresponding pair \(\{t,-t\}\); quotienting by
\(\iota\) gives the five-point description of
\(\Gamma_{1,2}\). The low-genus comparison in
\cite[Lemma~17(b)]{H17} identifies the action there with
the standard braid action of \(g\), up to an inner
correction. On the pure five-point group its class is also
that of \(g'\). Faithfulness of the GT factor gives
\begin{equation}
                         g'=g.                             \label{eq:harmonic-beta-q3}
\end{equation}
Every labelled forgetting map used below therefore carries
\(\delta_{10}\) to the standard action of \(g\), up to a
pure inner automorphism. On a mixed subgroup this follows
from the pure assertion: two automorphisms that agree on
the open pure group agree, since its centralizer is trivial.

\subsection{The harmonic four-division construction}

Fix a nonzero \(t_0\in T[2]\), and restrict \(N_0\) to its
stabilizer. Modulo its central involution this is the modular
subgroup \(R_0\simeq C_2*\mathbf Z\) of index three.
In the ten-point quotient retain \(b_0\) and the following
four branch points for a new double cover: the images of
the two nonzero two-torsion points other than \(t_0\), and
the two pairs \(\{t,-t\}\) with \(2t=t_0\).
The stabilizer of \(t_0\) preserves this set and its two
pairs, so restriction followed by this forgetting map is
defined. The two points above \(b_0\) in the new double cover
are individually marked in this order.

Here is the entire algebraic identification of the resulting
map. Write an elliptic curve with its chosen two-torsion
point as
\[
 E:\ y^2=x(x^2+Ax+B),\qquad t_0=(0,0).
\]
The other two-torsion abscissae are the roots of
\(x^2+Ax+B\). The halves of \(t_0\) have abscissae
\(\pm\sqrt B\), as follows directly from the duplication
formula. Hence the new double cover is
\begin{equation}
 E':\ z^2=(x^2+Ax+B)(x^2-B)
              =x^4+Ax^3-ABx-B^2.                           \label{eq:harmonic-beta-q4}
\end{equation}
Put
\[
 X=2z+2x^2+Ax,\qquad
 Y=2Xx-\frac{A^2}2x+\frac A2X-AB.
\]
Expansion gives
\begin{equation}
 E':\quad Y^2=X^3-B(A^2-4B)X.                               \label{eq:harmonic-beta-q5}
\end{equation}
One point above infinity becomes the origin. The other is
\begin{equation}
 P=\left(\frac{A^2}4,\frac{A(8B-A^2)}8\right).               \label{eq:harmonic-beta-q6}
\end{equation}
Thus the new curve has constant \(j=1728\).
The two infinity branches are globally ordered: they are the
sections \(z/x^2=+1\) and \(z/x^2=-1\) of the monic quartic
\eqref{eq:harmonic-beta-q4}. Under every Weierstrass coordinate change preserving
the origin and \(t_0\), the coordinates scale as
\(x\mapsto c^2x\), \(y\mapsto c^3y\), and
\(z\mapsto c^4z\). The ratio \(z/x^2\) is unchanged.
Thus the two sections cannot be interchanged by monodromy;
the first gives \(O\), the second gives \(P\), and the
construction defines a map to the group of two individually
marked points \(\Gamma_{1,2}\).

For \(E_4:y^2=x^3-x\), its quotient by \(\mu_4\) has
coordinate \(w=x^2\), with orbifold orders four at \(0\),
two at \(1\), and a puncture at infinity. With
\(t=A^2/B\), equations \eqref{eq:harmonic-beta-q5}--\eqref{eq:harmonic-beta-q6} give
\begin{equation}
                      w=\frac{t^2}{16(t-4)}.                \label{eq:harmonic-beta-q7}
\end{equation}
This is a degree-two orbifold covering: it is branched at
\(t=0\) above the order-four point and at \(t=8\) above the
order-two point. Its two unramified cusps are \(t=4,\infty\).
The source therefore has one order-two point and two cusps,
as required for \(R_0\).
It follows, without using a congruence computation to
identify a profinite map, that the construction induces a
discrete injection
\[
 \varphi_4:R_0\hookrightarrow G_4,\qquad
 G_4=q^{-1}(\langle s\rangle)\subset\Gamma_{1,2},
 \quad s=uvu,
\]
whose image has index two. Here \(q\) forgets the second
point and \(s\) has order four in \(\Gamma_{1,1}\).
The group \(G_4\) is the orbifold group of
\((E_4\setminus\{O\})/\mu_4\), hence is \(C_4*C_2\).
These inclusions retain the full profinite topology. Indeed the
point-pushing kernel \(\ker q\) has the full induced topology and
index four in \(G_4\). The quotient \(G_4/\ker q=\langle s\rangle\)
is detected injectively by signed reduction modulo four. Given a
finite-index subgroup of \(G_4\), choose an ambient open subgroup
whose intersection with \(\ker q\) is contained in it, and intersect
with this congruence kernel. The resulting intersection with
\(G_4\) lies in the prescribed subgroup. This proves the assertion
for \(G_4\); its index-two subgroup \(\varphi_4(R_0)\) inherits it.

Choose the conjugate of \(R_0\) generated by \(x=u^2\) and
the order-two image of \(s\). The cusp \(x\) has width two
and is \(t=\infty\); the other cusp \(t=4\) has width one.
The rotation character induced by \(\varphi_4\) is
\begin{equation}
                        x\longmapsto i,\qquad
                        s\longmapsto-1,                    \label{eq:harmonic-beta-q8}
\end{equation}
or its complex conjugate. This also follows from the
ramification in \eqref{eq:harmonic-beta-q7}: the order-two point lifts the square
of the order-four stabilizer, while a positive cusp loop
has primitive fourth-root rotation.

\subsection{Comparison on integral homology}

The normalized harmonic identity gives
\[
 \beta_g(u^2,v^2)^{-1}=f_g(u^2,s).
\]
Consequently the restriction of \(\Phi_g\) to \(R_0\) is
exactly the standard action in the generators \eqref{eq:harmonic-beta-q8}:
\begin{equation}
 x\longmapsto x,\qquad
 s\longmapsto f_g(x,s)^{-1}s f_g(x,s).                       \label{eq:harmonic-beta-q9}
\end{equation}
Let \(\mathcal A_g\) be the braid action on
\(\widehat\Gamma_{1,2}\), let
\(\widetilde\beta_g=\beta_g(a^2,b^2)\), and put
\(\Psi_4=\operatorname{Ad}(\widetilde\beta_g^{-1})\mathcal A_g\).
The action on the chain subgroup is \(\Phi_g\), so
\(\Psi_4\) fixes \(aba\), projects to an action fixing \(s\),
and preserves \(G_4\).

Equations \eqref{eq:harmonic-beta-q2}--\eqref{eq:harmonic-beta-q3}, restriction to the stabilizer of \(t_0\),
and the specified forgetting map give
\begin{equation}
 \Psi_4\varphi_4
       =\operatorname{Ad}(h)\varphi_4\Phi_g|_{R_0}            \label{eq:harmonic-beta-q10}
\end{equation}
for an \(h\in\widehat\Gamma_{1,2}\). More explicitly, write the
descended target action as
\(\delta_5=\operatorname{Ad}(c_3)\mathcal A_g\), with \(c_3\) pure.
Restriction and forgetting give
\[
 \delta_5\varphi_4=
 \operatorname{Ad}(\varphi_4(c_2))\varphi_4\Phi_g|_{R_0},
 \qquad
 h=\widetilde\beta_g^{-1}c_3^{-1}\varphi_4(c_2).
\]
The source inner correction \(c_2\) really belongs to the
stabilizer of \(t_0\), because it is pure modulo two; thus
its image under \(\varphi_4\) is defined. The target correction
comes from the pure five-point group, so is also defined in
\(\widehat\Gamma_{1,2}\). These account for all inner
corrections in \eqref{eq:harmonic-beta-q10}.
The rotation character \eqref{eq:harmonic-beta-q8} is unchanged by \eqref{eq:harmonic-beta-q9}.
Projecting \eqref{eq:harmonic-beta-q10} by \(q\) therefore shows that \(q(h)\)
centralizes \(\langle s\rangle\). In the completed modular
group the centralizer of its order-two factor is that factor;
lifting to \(\widehat\Gamma_{1,1}\) gives
\(q(h)\in\langle s\rangle\). Hence \(h\in G_4\).

Let \(K=\ker q=\pi_1(E_4\setminus O)\).
As in the integral cubic comparison, \(\mathcal A_g\) acts
as the identity on \(\widehat K^{\rm ab}\) for \(g\in N\).
For clarity, its action commutes with the full signed
congruence image, hence is scalar, and it fixes the
primitive class of \(ca^{-1}\). The latter assertion follows
from
\[
 \mathcal A_g(c)=P_g^{-1}cP_g,\qquad
 P_g=f_g(ba^2b,c^2)\in\widehat K,\qquad [P_g,a]=1.
\]
Thus the scalar is one. The left side of \eqref{eq:harmonic-beta-q10} consequently
acts on this homology by \(q(\beta_g)^{-1}=W_g\).

On the source, the kernel of \eqref{eq:harmonic-beta-q8} is the fundamental group
of a twice-punctured torus. After filling its two punctures,
its integral homology is a square lattice. Its map to
\(K^{\rm ab}\) is the inclusion of an index-two square
lattice: filling the two unramified cusps in \eqref{eq:harmonic-beta-q7} gives an
unramified double cover of the compact torus.
The source action on this lattice is multiplication by
\(\nu_{4,g}=B_g(i,-1)\). Here is a direct calculation.
The quotient is the Euclidean triangle group obtained from
\(R_0\) by imposing \(x^4=(xs)^4=1\). These are the two
filled positive cusp loops, so the peripheral action \eqref{eq:harmonic-beta-q9}
preserves their normal closure and descends to this quotient.
In the affine square-lattice quotient write
\[
 x=(a,i),\qquad s=(b,-1),\qquad
 [x,s]=(2a+(i-1)b,1)=:(c,1).
\]
The scalar \(c\) is nonzero and generates a finite-index
ideal in the lattice. Since
\([f_g]=R_g(X,Y)[x,s]\), the image of \(f_g(x,s)\) is
\((R_g(i,-1)c,1)\). If the induced affine automorphism
acts on translations by \(\nu\), allow a change of origin
by \(d\) in the rational quadratic algebra at each prime.
Every crossed translation term has this form, since
\(1-i\) is invertible in that algebra. Fixing \(x\) gives
\((\nu-1)a=(i-1)d\); the prescribed conjugation of \(s\)
gives \((\nu-1)b+2d=-2R_g(i,-1)c\).
Eliminating \(d\), and cancelling the nonzero lattice
factor \(c\), yields
\begin{equation}
               \nu=1-2(i-1)R_g(i,-1)=B_g(i,-1).             \label{eq:harmonic-beta-q11}
\end{equation}
Cancellation is valid in every completed integral component;
the ideal inclusion is torsion-free, including at two.

The inner automorphism by \(h\in G_4\) acts on
\(K^{\rm ab}\) by a power \(S^j\) of the order-four rotation.
The integral map between the two square lattices
intertwines their rotations, possibly with complex
conjugation. Equation \eqref{eq:harmonic-beta-q10} gives
\begin{equation}
                     W_g=S^j\nu_{4,g}^{\,\pm},              \label{eq:harmonic-beta-q12}
\end{equation}
where the superscript records the one geometric choice of embedding of
\(\widehat O_4\), fixed for all \(g\) and all coefficient primes.

Finally both nonrotation factors are one modulo four.
This is established for \(W_g\). For the beta value, the
class-two three-cycle relation gives \(R_g(1,1)=0\).
In \(\mathbf Z_2[i]\), with \(\pi=i-1\), the differences
\(i-1\) and \(-1-1\) have valuations one and two.
Thus \(R_g(i,-1)\in\pi\mathbf Z_2[i]\), and \eqref{eq:harmonic-beta-q11} is
one modulo \(\pi^4=(4)\). Since \(S\) has order four
modulo four, \eqref{eq:harmonic-beta-q12} forces \(j=0\).
This proves the asserted integral comparison. Reduction,
the cyclic order of the norm-one residue group, and the
bijectivity of squaring on its pro-\(\ell\) kernel give
the stated equivalent finite test.

This comparison does not prove that the displayed beta
value is always a square on \(D\). It identifies the
harmonic residue with an exact finite beta value; a further
relation is still needed to settle its tame square class.
